\chapter{Development Guidelines}
\label{ch:development-guidelines}

\section{Purpose}

No formal style guide, linter configuration or \texttt{CONTRIBUTING} file exists
in the repository. This chapter records the \emph{implicit} conventions that the
existing code follows, so that new code can be consistent with it. Conventions
derived from the code are labelled INFERRED; those stated by the authors are
labelled IMPLEMENTED.

\section{Repository layout and tooling}

\begin{itemize}
  \item The repository is a pnpm workspace monorepo with packages under
        \texttt{packages/*} (\texttt{core}, \texttt{cli}, \texttt{registry},
        \texttt{vscode-extension}).
  \item Node $\geq$ 18 is required; CI uses Node 20; pnpm 9 is pinned in CI.
  \item Shared root scripts: \texttt{build}, \texttt{dev}, \texttt{test},
        \texttt{validate}, \texttt{build:registry}, \texttt{serve}
        (all recursive where applicable).
  \item Build output directories (\texttt{dist}, \texttt{out}) and the generated
        registry modules are excluded by \texttt{.gitignore}.
\end{itemize}

\section{TypeScript conventions}

The shared compiler configuration (\texttt{tsconfig.base.json}) is the primary
contract:

\begin{longtable}{p{4.4cm}p{2.6cm}p{7.2cm}}
\caption{TypeScript settings and their implications}\\
\toprule
\textbf{Setting} & \textbf{Value} & \textbf{Implication for contributors} \\
\midrule

\texttt{strict} & \texttt{true} & No implicit \texttt{any}; null checks enforced.
Code must be fully typed. \\
\texttt{target} / \texttt{module} & ES2020 / CommonJS & Use ES2020 syntax; output is
CommonJS for the Node packages. \\
\texttt{moduleResolution} & \texttt{node} & Classic Node resolution; import
workspace packages by name. \\
\texttt{esModuleInterop} & \texttt{true} & Default imports from CommonJS modules
are permitted. \\
\texttt{declaration} / \texttt{declarationMap} & \texttt{true} & Public
\texttt{core} API is published with types. \\
\texttt{resolveJsonModule} & \texttt{true} & JSON can be imported as a module. \\
\bottomrule
\end{longtable}

\textbf{INFERRED coding style:} interfaces and type aliases are preferred for
public shapes (\texttt{ScreenRecipe}, \texttt{GenerationOptions},
\texttt{DetectionResult}); functions are small and single-purpose; utilities live
in a \texttt{utils/} folder within a package; the core exposes a small public
surface via \texttt{index.ts}. No barrel files beyond the package entry points
are used.

\section{Extension conventions (INFERRED)}

\begin{itemize}
  \item The extension is bundled with esbuild to CommonJS, targeting Node 18,
        with \texttt{vscode} marked external and source maps enabled.
  \item Runtime libraries the extension needs are bundled, so they are declared
        as \texttt{devDependencies} rather than \texttt{dependencies}.
  \item The extension entry point is \texttt{src/extension.ts}; commands,
        completion provider, generation and project logic are separated into
        their own modules.
  \item The offline catalog is regenerated at build time by
        \texttt{scripts/prepare-registry.mjs}; the generated file is never edited
        by hand.
\end{itemize}

\section{Recipe and template conventions (INFERRED + IMPLEMENTED)}

Recipes are the project's main "content" authoring surface, and the conventions
here are strong and consistent.

\subsection*{Recipe identity}

\begin{itemize}
  \item An id is \texttt{<category>/<name>}, lowercase, hyphen-separated
        (\texttt{\^{}[a-z0-9-]+/[a-z0-9-]+}) --- for example \texttt{auth/login}.
  \item The recipe file must live at a path that matches the id:
        \texttt{packages/registry/data/auth/login.json}. The validator enforces
        this.
  \item The recipe \texttt{version} is semver of the recipe itself
        (\texttt{auth/login} is at 0.3.0).
\end{itemize}

\subsection*{Recipe content}

\begin{itemize}
  \item \texttt{supported} must list only options the recipe actually
        implements; the clients reject unsupported combinations.
  \item Every output file is a separate entry in \texttt{files}, with a target
        \texttt{path} and a \texttt{templateFile}. Multiple entries can target the
        same path under different \texttt{when} conditions (for example TypeScript
        vs JavaScript, or native vs shadcn).
  \item \texttt{when} selects a file entry by styling, router, ui, icons and/or
        TypeScript.
  \item \texttt{variants} at the recipe level names alternative styles; a file
        maps a variant to its template via \texttt{files[].variants}.
  \item \texttt{primitives} lists UI primitive recipe ids required when
        \texttt{ui=shadcn} (for example \texttt{ui/utils}, \texttt{ui/button},
        \texttt{ui/input}).
  \item \texttt{dependencies} must include every third-party package the
        templates import. (The \texttt{glass} variants currently violate this;
        see Chapter~\ref{ch:known-issues}.)
\end{itemize}

\subsection*{Template naming}

Templates live under \texttt{packages/registry/templates/<category>/<name>/} and
follow a dotted convention that mirrors the output location, with optional
modifier segments before the extension:

\begin{longtable}{p{6cm}p{8.2cm}}
\caption{Template naming examples}\\
\toprule
\textbf{Template file} & \textbf{Meaning} \\
\midrule

\texttt{auth/login/app.login.index.hbs} & Expo Router \texttt{index} screen
(default). \\
\texttt{auth/login/app.login.index.shadcn.hbs} & shadcn variant of the same
screen. \\
\texttt{auth/login/app.login.index.glass.hbs} & \texttt{glass} variant of the same
screen. \\
\texttt{auth/login/components.login-form.hbs} & A component file. \\
\texttt{auth/login/hooks.use-login.hbs} & A hook file (shared by TS and JS
outputs). \\
\texttt{auth/login/screens.loginscreen.react-navigation.hbs} &
\texttt{react-navigation} screen variant. \\
\bottomrule
\end{longtable}

\subsection*{Template content rules (INFERRED)}

\begin{itemize}
  \item Templates are TypeScript/JavaScript source, not markup; they render to
        \texttt{.tsx}, \texttt{.ts} or \texttt{.js} files.
  \item Handlebars is used only with the two registered helpers
        (\texttt{eq}, \texttt{neq}); contributors should not assume custom helpers
        exist.
  \item New recipe files must be referenced by path from a recipe document;
        an orphaned \texttt{.hbs} file is not shipped.
\end{itemize}

\section{Versioning and change control (INFERRED)}

\begin{itemize}
  \item The public packages are versioned in their \texttt{package.json} and
        released by pushing a \texttt{v*} tag.
  \item The extension follows the CHANGELOG convention (Keep a Changelog plus
        Semantic Versioning), as stated at the top of its \texttt{CHANGELOG.md}.
  \item Recipe changes are versioned per recipe, independently of the packages.
\end{itemize}

\section{Recommended additions to these guidelines}

\textbf{Status:} RECOMMENDED.

\begin{itemize}
  \item Add a linter and formatter (for example ESLint and Prettier) with a
        committed configuration, and run them in CI.
  \item Add an \texttt{.editorconfig} to fix indentation and line endings across
        editors.
  \item Document the separation of concerns explicitly: rendering belongs in
        \texttt{core}; orchestration should also be shared (see
        Chapter~\ref{ch:technical-debt}); clients should hold only
        environment-specific code.
  \item Require a recipe to declare every package its templates import, with an
        automated check.
\end{itemize}

\section{Summary}

\textbf{Status:} INFERRED (most conventions) / IMPLEMENTED (recipe identity,
schema, TypeScript settings, extension bundling). The codebase is internally
consistent and easy to imitate; the gap is that none of this is written down and
none of it is enforced automatically.
